\subsection{Scheduling procedure}
\label{subsec:scheduling_procedure}
\label{subsec:scheduling}

We represent the chamber settings returned by
\(\mathrm{ChamberAssignment}\) by \(\mathbf{x}\) and the
sample-size vector by \(\mathbf{s}\).
An integral component of the proposed solution approach is to
evaluate the quality of a solution.
We need the total tardiness of a given pair
\((\mathbf{x},\mathbf{s})\).
For that purpose we use an earliest-due-date (EDD) constructive
scheduler, called \(\mathrm{Schedule}\) in the outer framework,
that assigns tests to stations.
Such constructive heuristics are particularly attractive for
unrelated parallel-machine problems with eligibility constraints
and due-date pressure, where repeated implementation must remain
computationally efficient
\citep{PerezGonzalezFernandezViagasZamoraGarciaFraminan2019,WangHeKerkhoveVanhoucke2017}.

Readiness follows the stage logic of the problem definition.
Gas amount determination, when required, is executed once per
product and may start at time zero.
Pull-down is executed once per sample and can start only after
gas amount determination of that product has finished.
Third-stage tests are released after every pull-down of that
product has finished.
A consumer-usage test is released after every pull-down of that
product has finished and every required third-stage test has
started.

At each iteration the procedure rebuilds a candidate pool
\(\mathcal{O}^{\mathrm{cand}}\) with at most one operation per
product.
That operation is returned by \(\mathrm{BestReady}(p)\).
The inputs of \(\mathrm{BestReady}(p)\) are the required tests
\(\mathcal{T}_p\) of product \(p\) together with their durations
and environments, the voltage requirement of \(p\), the current
stage state of \(p\), the \(s_p\) sample-availability clocks
induced by \(\mathbf{s}\), and the stations made eligible by
\(\mathbf{x}\).
It does not read the full schedule \(\Pi\).
It inspects only the currently released tests of \(p\).
Let \(\rho_{pt}\) denote the stage-release time of test \(t\) of
product \(p\), let \(\alpha_{p\ell}\) denote the availability time
of sample \(\ell\) of \(p\), and let \(\beta_{cr}\) denote the
availability time of station \(r\) in chamber \(c\).
For each released test, \(\mathrm{BestReady}\) enumerates
compatible stations and samples and records the earliest feasible
start
\[
\max\{\rho_{pt},\alpha_{p\ell},\beta_{cr}\}.
\]
A station is compatible if its chamber matches the test
environment, is humidity-capable when the test is
humidity-controlled, and is voltage-capable when the product
requires it.
Among those released tests it returns the operation with the
earliest planned start, already paired with that station and
sample, or reports that \(p\) has nothing ready.

A candidate in \(\mathcal{O}^{\mathrm{cand}}\) is therefore not an
unbound test waiting only on a sample and a station.
Stage release is required before the test can enter the pool, and
that release time is included in the planned start.
The candidate already carries one compatible station and one
sample; its start is the first time at which the stage, the
sample, and the station are all free.
While it sits in the pool it waits only to be selected.
The next operation in the system is the candidate with the
earliest due date; ties are broken by the earliest feasible start.
Only that winner is written onto its planned station and sample.
Station and sample clocks, and the stage state of its product,
are then updated, and the pool is rebuilt from the new state.
Unselected candidates occupy no resource.
When no product has a ready operation, the completion time
\(C_p\) is the finish time of product \(p\)'s last scheduled
test, and \(T^{tot}\) is computed as in the problem definition.
When a full schedule is stored, an independent check audits
required-test multiplicity, environment, humidity and voltage
eligibility, station non-overlap, sample non-overlap, and stage
release.
Algorithm~\ref{alg:schedule} is the procedure.

\begin{algorithm}[H]
\caption{Schedule}
\label{alg:schedule}
\label{alg:scheduling}
\label{alg:scheduling_procedure}
\label{alg:scheduler_short}
\begin{algorithmic}[1]
\State \textbf{Input:} instance data, chamber assignment \(\mathbf{x}\), sample vector \(\mathbf{s}\)
\State \textbf{Output:} schedule \(\Pi\), total tardiness \(T^{tot}\)
\State \(\Pi\leftarrow\emptyset\); initialize station availability, sample availability, and stage states \label{step:sch-init}
\Repeat \label{step:sch-loop}
    \State \(\mathcal{O}^{\mathrm{cand}}\leftarrow\emptyset\) \label{step:sch-pool}
    \For{each product \(p\in\mathcal{P}\)}
        \State \(o_p\leftarrow\mathrm{BestReady}(p)\) \label{step:sch-ready}
        \If{\(o_p\) exists}
            \State \(\mathcal{O}^{\mathrm{cand}}\leftarrow\mathcal{O}^{\mathrm{cand}}\cup\{o_p\}\)
        \EndIf
    \EndFor
    \If{\(\mathcal{O}^{\mathrm{cand}}\neq\emptyset\)} \label{step:sch-select}
        \State Choose \(o^{\star}\in\mathcal{O}^{\mathrm{cand}}\) with the earliest due date; break ties by the earliest feasible start
        \State Assign \(o^{\star}\) to its planned compatible station and sample \label{step:sch-assign}
        \State Update station availability, sample availability, and the stage state of its product \label{step:sch-update}
        \State Append \(o^{\star}\) to \(\Pi\)
    \EndIf
\Until{\(\mathcal{O}^{\mathrm{cand}}=\emptyset\)} \label{step:sch-stop}
\State Compute \(C_p\), \(T_p\), and \(T^{tot}\) \label{step:sch-obj}
\State \Return \((\Pi,T^{tot})\)
\end{algorithmic}
\end{algorithm}

Step~\ref{step:sch-init} starts from an empty schedule and zero
availability clocks.
Each pass (Step~\ref{step:sch-loop}) rebuilds the candidate pool
(Step~\ref{step:sch-pool}).
Step~\ref{step:sch-ready} returns, for each product, at most one
released operation, already bound to a station and a sample, whose
planned start is \(\max\{\rho_{pt},\alpha_{p\ell},\beta_{cr}\}\).
Step~\ref{step:sch-select} applies EDD to the pool, and
Step~\ref{step:sch-assign} writes the winner onto the station and
sample already chosen by \(\mathrm{BestReady}\).
Step~\ref{step:sch-update} advances the clocks so that later
passes see the new stage state.
The remaining candidates of that pass are discarded.
The loop stops when no product has a ready operation
(Step~\ref{step:sch-stop}).
Step~\ref{step:sch-obj} records tardiness.
